% Generated draft: requires mathematical and source review.
\documentclass[11pt]{article}
\usepackage[margin=25mm]{geometry}
\usepackage{fontspec}
\setmainfont{Latin Modern Roman}
\usepackage{amsmath,amssymb,mathtools,unicode-math}
\setmathfont{Latin Modern Math}
\usepackage{enumitem,needspace,xurl,hyperref}
\setlength{\parindent}{0pt}
\setlength{\parskip}{5pt}
\setlength{\emergencystretch}{3em}
\newcommand{\N}{\mathbb{N}}
\newcommand{\Z}{\mathbb{Z}}
\newcommand{\Q}{\mathbb{Q}}
\newcommand{\R}{\mathbb{R}}
\newcommand{\C}{\mathbb{C}}
\newcommand{\F}{\mathbb{F}}
\newcommand{\A}{\mathbb{A}}
\newcommand{\PP}{\mathbb{P}}
\newcommand{\E}{\mathbb{E}}
\newcommand{\eps}{\varepsilon}
\newcommand{\dd}{\,\mathrm{d}}
\begin{document}
\section*{UnsolvedMath ID: 11100053}
\textbf{Source problem code:} catalogue entry
\subsection{Definitions and mathematical statement}
Let $A$ be a cofibrant operad, as in the source. The ambient category, model structure on operads, and ``above results'' to which the item refers are not included in the supplied record. Subject to specifying an appropriate homotopical setting, the problem asks for two constructions:
\begin{enumerate}
\item For a given object $X$, construct spectral sequences converging to the homotopy groups of the space of $A$-algebra structures on $X$.
\item For two given $A$-algebras $B$ and $C$, construct spectral sequences converging to the homotopy groups of the space of $A$-algebra maps from $B$ to $C$.
\end{enumerate}
An $A$-algebra structure means a compatible collection of operations parametrized by $A$. In a setting with endomorphism operads and suitable derived mapping spaces, a possible interpretation of the first target is
\[
 \mathcal S_A(X)=\mathbf R\operatorname{Map}_{\mathrm{Operad}}
       (A,\operatorname{End}_X),
\]
where the arity-$r$ part of $\operatorname{End}_X$ represents maps $X^{\otimes r}\to X$, with the appropriate enriched and derived replacements. The second target is
\[
 \mathbf R\operatorname{Map}_{\operatorname{Alg}_A}(B,C).
\]
These expressions are conditional interpretations: the source does not specify whether the first space fixes the underlying object with a chosen identification or quotients by equivalences, nor the model structures ensuring the displayed derived mapping spaces exist with the intended meaning.

For each relevant basepoint and each homotopy degree for which they are defined, the requested spectral sequences should recover the homotopy groups of these targets under explicit convergence hypotheses. The source gives no pages, bidegrees, differentials, filtrations, or connectivity/completeness assumptions, so these cannot be filled in from this record alone. It singles out the $A_{\infty}$ operad: the corresponding spectral sequences are described as key formal ingredients in the Hopkins--Miller proof of an action of the stabilizer group on Morava $E$-theory. The height, prime, group convention, and coherence of that action are not specified here.
\subsection{Short English statement}
Construct convergent spectral sequences for spaces of algebra structures over a cofibrant operad and for mapping spaces between its algebras. Preserve the stated $A_{\infty}$ and Morava $E$-theory motivation. The registry reports this solved by operadic deformation theory, but the omitted antecedents and unspecified hypotheses prevent a complete theorem statement from this isolated item.
\subsection{Sources}
\begin{itemize}
\item Local UnsolvedMath registry, permanent ID $11100053$, source code \texttt{catalogue entry}, proposed by Mark Hovey; created 2026-07-29, updated 2026-08-06. Supplied status \texttt{solved}; research classification \texttt{SOLVED-IN-LITERATURE}.
\item Mark Hovey, \emph{Hovey's Algebraic Topology Problem List}, ``Model categories,'' item 7, surviving mirror \url{https://www-users.cse.umn.edu/~tlawson/hovey/}. Accessed 2026-07-29; HTML ordered-list extraction; 65 top-level items in nine category pages; default difficulty L3.
\item The statement attributes the Morava $E$-theory application to Hopkins--Miller. The research summary names Rezk, Hinich, and Goerss--Hopkins operadic deformation theory and says that it realizes ``Hovey's step 4'' and the Morava-$E$ application. It supplies no exact references or statement of that step; the extracted source item is numbered 7. No identification of these labels is invented here.
\item The background separately marks source-status evidence \texttt{NEEDS\_REVIEW}: the source presents an open problem and a current primary-literature audit remains. Source-rights note \texttt{NEEDS\_REVIEW}: publicly accessible author/publisher copy, redistribution terms requiring review. Dataset attribution: Ulam AI / UnsolvedMath Contributors, CC BY 4.0.
\end{itemize}
\end{document}
